\usetikzlibrary{shadows.blur, fit, backgrounds}  % or shadows

\resizebox{\linewidth}{!}{%

\begin{tikzpicture}[
    % ---- visual style ------------------------------------------------
    bar/.style = {
    draw=black, line width=0.4pt,
    minimum height=14mm,
    inner sep=4pt, outer sep=0pt,
    anchor=west,
    font=\sffamily\itshape\small,
    blur shadow={shadow blur steps=4, shadow blur extra rounding=0.5pt,
                 shadow xshift=1pt, shadow yshift=-1pt, shadow opacity=25}
    },
    warmup/.style = {bar, fill=warmupgreen},        % light blue (like "Surrogate
    hybrid/.style = {bar, fill=hybridpurple},
    node/.style   = {bar, fill=nodegrey},          % soft pink (like "Primal Solver")
    stage/.style  = {bar, fill=yellow!40!white,        % warm yellow (like "Adjoint     Solver")
                     text=black},
    empty/.style  = {bar, fill=white},
    tick/.style      = {font=\sffamily\footnotesize},
    cond/.style      = {font=\sffamily\footnotesize, align=center},
    arr/.style       = {-{Latex[length=2mm]}, thin},
    % ---- block lengths (tweak freely) --------------------------------
    /lengths/.cd,
      warmup/.initial    = 35mm,
      empty/.initial     =20mm,
      gap1/.initial      = 5mm,    % small solver run before pipeline starts
      emptygap/.initial = 25mm,
      solverA/.initial   = 12mm,    % solver before first NODE rollout
      nodeA/.initial     = 8mm,
      solverB/.initial   = 8mm,
      nodeB/.initial     = 22mm,    % main NODE rollout block
      solverC/.initial   = 8mm,
      nodeC/.initial     = 6mm,    % "Updating data" pink stretch
      retrain/.initial   = 10mm,    % tail solver run before retraining
      gap2/.initial      = 12mm,    % short solver after retraining
      nodeD/.initial     = 10mm,
      tail/.initial      = 6mm,
]

% --- shorthand to read the lengths -----------------------------------
\def\L#1{\pgfkeysvalueof{/lengths/#1}}

% =====================================================================
% TIMELINE BAR (laid out left-to-right by chaining anchors)
% =====================================================================
\node[warmup, minimum width=\L{warmup}]                         (warmup)  at (0,0) {Initial data-collection};

\node[stage, minimum width=\L{emptygap}, right=0pt of warmup] (train1) {$\mathcal{P}$ training};

\node[node,   minimum width=\L{nodeA},   right=0pt of train1]   (nA) {};
\node[hybrid, minimum width=\L{solverB}, right=0pt of nA]       (sB) {};
\node[node,   minimum width=\L{nodeB},   right=0pt of sB]       (nB) {NODE rollout};
\node[hybrid, minimum width=\L{solverC}, right=0pt of nB]       (sC) {};
\node[node,   minimum width=\L{nodeC},   right=0pt of sC]       (nC) {};

\node[warmup, minimum width=\L{retrain}, right=0pt of nC]       (sD) {Updating data};
\node[stage, minimum width=\L{emptygap}, right=0pt of sD]       (train2) {$\mathcal{P}$ retraining};


\node[node,   minimum width=\L{nodeA},   right=0pt of train2]   (nD) {};
\node[hybrid, minimum width=\L{solverB}, right=0pt of nD]       (sF) {};
\node[node,   minimum width=\L{nodeB},   right=0pt of sF]       (nE) {NODE rollout};
\node[hybrid, minimum width=\L{solverC}, right=0pt of nE]       (sG) {};
\node[node,   minimum width=\L{nodeC},   right=0pt of sG]       (nF) {};

\node[warmup, minimum width=\L{tail},    right=0pt of nF, text width=20mm]      (tail) {\scriptsize Hybrid-Retrain loop continues};

% =====================================================================
% TIME AXIS (drawn just below the bar)
% =====================================================================
\coordinate (axL) at ($(warmup.south west)+(0,-1mm)$);
\coordinate (axR) at ($(tail.south east)+(4mm,-1mm)$);
\draw[arr, thick] (axL) -- (axR);

% --- tick labels on the axis -----------------------------------------
\node[tick, below=2mm of warmup.south east,anchor=north] {$t_{\text{train}}$};
\node[tick, below=2mm of train1.south east,anchor=north] {$t_{\text{handoff}}$};

\node[tick, below=2mm of nC.south east,    anchor=north] {$n_\text{flags} \geq n_\text{max, flags}$};
\node[tick, below=2mm of nF.south east,    anchor=north] {$n_\text{flags} \geq n_\text{max, flags}$};
\node[tick, below=2mm of tail.south east,  anchor=north] {$t_{\text{end}}$};

% =====================================================================
% OOD-trigger condition d(k,z) > tau, pointing at NODE-rollout boundaries
% =====================================================================
\node[cond] (ood) at ($(nB.north)+(0,7mm)$) {$d^1_k(\tilde{z}; \mathcal{K}_1) > \tau$};
\foreach \t in {nA.north east, nB.north east, nC.north east}{%
  \draw[arr] (ood) -- (\t);
}

\node[cond] (ood2) at ($(nE.north)+(0,7mm)$) {$d^2_k(\tilde{z}; \mathcal{K}_2) > \tau$};
\foreach \t in {nD.north east, nE.north east, nF.north east}{%
  \draw[arr] (ood2) -- (\t);
}

\node[cond] (switch1) at ($(nB.south)-(0,7mm)$) {$n_{step} \bmod n_{switch} = 0 $};
\foreach \t in {nA.south west, nB.south west, nC.south west}{%
  \draw[arr] (switch1) -- (\t);
}

\node[cond] (switch2) at ($(nE.south)-(0,7mm)$) {$n_{step} \bmod n_{switch} = 0 $};
\foreach \t in {nD.south west, nE.south west, nF.south west}{%
  \draw[arr] (switch2) -- (\t);
}

\begin{scope}[on background layer]
  % --- 1. Outer container first ---------------------------------------
  \node[draw=black!50, dashed, rounded corners=4pt, fill=none, inner sep=6mm, draw=none, fit=(warmup)(tail)(ood)] (outer) {};

  % --- 2. Shared top/bottom Y reference (just points; we'll project) --
  \coordinate (topY) at ($(outer.north) + (0,-2mm)$);
  \coordinate (botY) at ($(outer.south) + (0, 2mm)$);

  % --- 3. Phase panel style -------------------------------------------
  \tikzset{
    phase/.style={
      draw=black!40, dashed, rounded corners=3pt,
      inner xsep=1mm, inner ysep=0pt
    }
  }

  % --- 4. Phase panels: project topY/botY onto each panel's own x range
  \node[phase, fill=phasewarmup,   fit=(warmup) (warmup.west |- topY) (warmup.east |- botY)]              (pW)  {};
  \node[phase, fill=phasetraining, fit=(train1) (train1.west |- topY)   (train1.east  |- botY)]     (pT)  {};
  \node[phase, fill=phasehybrid,   fit=(nA)(sB)(nB)(sC)(nC)(ood) (nA.west |- topY)   (nC.east  |- botY)]  (pH1) {};
  \node[phase, fill=phaseretrain,  fit=(sD)(train2)(sD.west |- topY)   (train2.east  |- botY)]      (pR)  {};
  \node[phase, fill=phasehybrid,   fit=(nD)(sF)(nE)(sG)(nF)(ood2)(nD.west |- topY)   (nF.east  |- botY)]                    (pH2) {};
  \node[phase, fill=phaseretrain,  fit=(tail)(tail.west |- topY)   (tail.east  |- botY)]                  (pR2) {};
\end{scope}

% \node[node,   minimum width=\L{nodeA},   right=0pt of train2]   (nD) {};
% \node[hybrid, minimum width=\L{solverB}, right=0pt of nD]       (sF) {};
% \node[node,   minimum width=\L{nodeB},   right=0pt of sF]       (nE) {NODE rollout};
% \node[hybrid, minimum width=\L{solverC}, right=0pt of nE]       (sG) {};
% \node[node,   minimum width=\L{nodeC},   right=0pt of sG]       (nF) {};

\foreach \panel/\txt in {%
  pW/Warmup,
  pT/Training,
  pH1/Hybrid loop,
  pR/Retraining,
  pH2/Hybrid loop,
  pR2/Retraining%
}{%
  \node[font=\sffamily\itshape\footnotesize, anchor=south, yshift=+1.5mm]
    at (\panel.north) {\txt};
}

\end{tikzpicture}}
